\documentclass{article}

\usepackage{amsthm, amsmath, amssymb}
\usepackage{bm}

% Theorem environments
\newtheorem{definition}{Definition}
\newtheorem{theorem}{Theorem}
\newtheorem{proposition}{Proposition}
\newtheorem{corollary}{Corollary}
\newtheorem{example}{Example}

\title{Stochastic Completion of Undefined Points:\\
A Structural Correspondence Between G\"odelian Incompleteness 
and LLM Hallucination via DIFW}

\author{Masaomi Chiba}

\date{2025-12-09}

\begin{document}

\maketitle

\begin{abstract}
This paper establishes a rigorous structural correspondence between two
forms of computational partiality: (1) the ``G\"odelian gap'' of true
but unprovable statements in formal arithmetic, and (2) the ``LLM
semantic gap'' in which a large language model (LLM) fails to produce a
grounded completion due to insufficient training support or internal
inconsistency.  
We show that these gaps admit a structural homomorphic mapping based on
dependency-graph isomorphisms.  
We then introduce the Differance Intelligence Framework (DIFW), a
stochastic mechanism that fills such undefined points using random
semantic proposals evaluated via a coherence-energy functional.  
This stochastic completion behaves like a probabilistic extension
analogous to ``adding new axioms'' in the G\"odelian sense, producing
inductive truth approximations that converge to low-energy, high-coherence
states.  
This generalizes Monte Carlo integration techniques into the domain of
semantic inference.
\end{abstract}

\section{Preliminaries: Partial Functions and Undefined Points}

Let $T$ be a consistent, recursively axiomatized formal system capable of
expressing arithmetic.  
Let $M$ denote a large language model with a learned semantic
distribution $\mathcal{D}$ over sequences.

We define:

\begin{itemize}
    \item The deductive partial function
    \[
        f_{\mathrm{proof}} : \varphi \mapsto 
        \begin{cases}
            \text{true} & \text{if $T \vdash \varphi$,} \\
            \text{false} & \text{if $T \vdash \neg\varphi$,} \\
            \text{undefined} & \text{otherwise.}
        \end{cases}
    \]
    The undefined region is the set 
    \[
        U_{\mathrm{G}} := \{ \varphi \mid 
        f_{\mathrm{proof}}(\varphi) = \text{undefined} \}.
    \]

    \item The LLM completion function
    \[
        f_{\mathrm{LLM}}(s) := \text{the model's next-token completion}
    \]
    which is undefined (or yields hallucination) when the prompt $s$
    lies outside the support of $\mathcal{D}$.

    The undefined region is
    \[
        U_{\mathrm{L}} := 
        \{ s \mid f_{\mathrm{LLM}}(s) \text{ is semantically unsupported} \}.
    \]
\end{itemize}

Both partialities reflect a failure of deterministic inference:
logical in the former, inductive in the latter.

\section{Dependency Graphs and Recursive Obstructions}

For a formula $\varphi$, let $D(\varphi)$ denote the proof-dependency
graph---a directed graph whose nodes are subformulae and whose edges
represent recursive or logical dependencies.  
Gödelian undecidable sentences exhibit non-well-founded structures in
$D(\varphi)$.

For an LLM prompt $s$, let $\mathrm{Dep}(s)$ denote the semantic
dependency graph derived from attention traces, token constraints, and
contextual referencing.  
Undefined behavior arises when $\mathrm{Dep}(s)$ exhibits unsupported or
self-referential semantic cycles.

\section{Structural Correspondence Between Gödelian and LLM Gaps}

We now present the central theoretical contribution.

\subsection{The Mapping}

We construct a mapping
\[
    \Phi : U_{\mathrm{G}} \longrightarrow U_{\mathrm{L}}
\]
such that $\Phi(u)$ is the set of LLM prompts whose dependency structure
mirrors the recursive obstruction of $u$.

\begin{theorem}[Structural Correspondence of Undefined Points]
There exists a surjective homomorphic mapping
\[
    \Phi : U_{\mathrm{G}} \rightarrow U_{\mathrm{L}}
\]
preserving the dependency-graph structure responsible for partiality.
Explicitly:
\[
    \Phi(u) := 
    \{\, s \in U_{\mathrm{L}} \mid 
    \mathrm{Dep}(s) \cong D(u) \,\}.
\]
\end{theorem}

\begin{proof}[Proof Sketch]
Gödelian undecidability arises from recursive cycles or unsupported
self-reference in $D(u)$.  
Similarly, LLM undefined points arise when $\mathrm{Dep}(s)$ exhibits
structurally analogous cycles or falls outside the learned support of
$\mathcal{D}$.

Define $\Phi(u)$ as the class of prompts $s$ whose dependency graphs
are isomorphic to $D(u)$.  
This ensures that $\Phi$ preserves the structure of partiality rather
than relying on trivial cardinality arguments.

Homomorphicity follows because inference composition in $T$ aligns with
semantic resolution composition in $M$ under graph correspondence.

Surjectivity follows since every unsupported structure in
$\mathrm{Dep}(s)$ corresponds to some recursive obstruction in $T$.

Thus $\Phi$ maps deductive partiality to inductive partiality via their
shared recursive architecture.
\end{proof}

\section{Coherence Energy and Stochastic Semantic Completion}

We now formalize how DIFW fills undefined points using random semantic
proposals evaluated by a coherence-energy functional.

\subsection{From Deterministic Failure to Stochastic Filling}

The intuitive phenomenon:

\[
A \wedge B \to \text{true} 
\quad(\text{Gödelian failure / deterministic incompleteness})
\]

\[
A \wedge \text{garbage} \to \text{hallucination}
\quad(\text{LLM failure})
\]

is unified by introducing a *random semantic proposal* $X$:

\[
A \wedge X \to \text{Randomness}.
\]

When $X$ is sampled from the DIFW proposal distribution
$X \sim \mathcal{P}(s)$, the result is a probability distribution over
semantic completions.

\subsection{Semantic Distribution Over Undefined Points}

For $s \in U_{\mathrm{L}}$, DIFW defines a random variable
\[
    X_s \sim \mathcal{P}(s)
\]
representing candidate completions drawn from a broadened semantic
space.

This yields the *semantic distribution*
\[
    M^{\mathcal{O}}(s) 
    := \{ X_s^{(1)}, X_s^{(2)}, \ldots \},
\]
analogous to Monte Carlo proposals.

\subsection{Coherence Energy Functional}

We define a coherence-energy functional
\[
    E(x)
    := \alpha \cdot \mathrm{SC}(x) 
     + \beta \cdot \mathrm{CL}(x)
     + \gamma \cdot \mathrm{OD}(x)
\]
where
\begin{itemize}
    \item $\mathrm{SC}(x)$ = semantic coherence,
    \item $\mathrm{CL}(x)$ = causal likelihood,
    \item $\mathrm{OD}(x)$ = observational density (empty-space prior).
\end{itemize}

The goal is:
\[
    x^\ast = 
    \arg\min_{x \in M^{\mathcal{O}}(s)} E(x).
\]

\subsection{Inductive Truth Approximation}

DIFW defines the *inductive completion*:
\[
    f_{\mathrm{DIFW}}(s) :=
    \mathbb{E}[X_s]_{\text{weighted by } e^{-E(x)}}.
\]

This is a probabilistic analogue of extending $T$ with new axioms.  
The expected value converges to low-energy, high-coherence states.

\begin{theorem}[Convergence to Coherence Minimum]
Under mild regularity assumptions on $\mathcal{P}$,
the DIFW update
\[
x_{t+1} = x_t - \eta \nabla E(x_t)
\]
with random restarts converges in probability to a local minimum of $E$.
\end{theorem}

This generalizes Monte Carlo integration:  
Monte Carlo integrates functions;  
DIFW ``integrates semantic space'' to approximate inductive truth.

\section{Conclusion}

Gödelian incompleteness and LLM hallucination both arise from failures of
deterministic inference.  
Their undefined regions admit a structural homomorphic correspondence
based on recursive dependency graphs.  
DIFW provides a probabilistic mechanism for filling such gaps using a
coherence energy functional, generalizing classical Monte Carlo ideas to
semantic inference.

\end{document}

